% ============================================================================
%  Related work -- % ============================================================================
%  Related work -- % ============================================================================
%  Related work -- % ============================================================================
%  Related work -- \input{sections/related}
% ============================================================================
\section{Related Work}
\label{sec:related}

\paragraph{Fish feeding intensity assessment.}
Early FFIA relied on observers or on hand-crafted image analysis, such as background segmentation and
the tracking of collective behaviour~\cite{zhao2017appetite}. These methods were sensitive to water
surface fluctuations and computationally heavy. Deep learning moved the field to convolutional
classifiers of feeding images~\cite{zhou2019ffia}, with later refinements such as dual attention on
EfficientNet-B2~\cite{yang2021effnetb2} and multi-scale fusion on MobileNetV3~\cite{zhang2023msif}.
Single images miss the dynamics of feeding, so video methods convert RGB sequences to optical flow and
classify them with 3-D CNNs~\cite{ubina2021ffia}, at a higher computational cost. Acoustic FFIA was
introduced with the AFFIA3K dataset, where log-mel spectrograms are classified by a CNN~\cite{cui2022affia3k}.
Surveys of fish behaviour analysis~\cite{an2021behaviour,yang2021smartfish} describe both modalities as
complementary but individually fragile. The audio-visual setting was established by AV-FFIA and
U-FFIA~\cite{cui2024uffia}. That work benchmarked audio models (PANNs~\cite{kong2020panns}, AST~\cite{gong2021ast}),
video models (I3D, S3D~\cite{xie2018s3d}, ViViT), and fusion by self-attention, cross-attention and
attention bottlenecks. It then proposed a unified mixed-modality model with modality dropout and
distillation from CNN10 and S3D teachers. Later works raised the reported AV accuracy with
cross-attention fusion~\cite{acafnet2025}, dynamic gating between experts~\cite{dynagate2026},
hierarchical multimodal fusion~\cite{hmmfnet2026}, heterogeneous distillation~\cite{hmkd2026} and
lightweight dual-modal designs~\cite{ldmf2026}. Each of them uses its own protocol: the official split,
cross-validation or a filtered subset. We report results under several protocols, including a
session-disjoint split, so that comparisons with these works are explicit.

\paragraph{Audio-visual fusion.}
Multimodal fusion is commonly divided into early, late and intermediate fusion~\cite{baltrusaitis2019mmml}.
Transformer fusion concatenates the token sequences, or uses cross-modal attention to let one modality
query the other. Attention bottlenecks restrict cross-modal flow to a few shared tokens~\cite{nagrani2021mbt},
which lowers cost and forces the model to condense information. Audio can also be used to make video
processing cheaper by previewing which frames matter~\cite{gao2020listen}. Self-supervised
audio-visual pretraining, such as contrastive masked autoencoders~\cite{gong2023cavmae}, learns joint
representations from large unlabelled corpora. Modality dropout~\cite{neverova2016moddrop} makes a fused model
robust to missing inputs and is used by U-FFIA to train a single model for audio, video or both. HCF
keeps the shared bottleneck tokens of~\cite{nagrani2021mbt}, but differs in two ways: the bottleneck
reads both modalities \emph{in parallel}, and their views and interactions are combined by an
algebraic product rather than by sequential attention or concatenation. We further attach HCF at any
subset of depths, with feedback into both streams.

\paragraph{Hypercomplex neural networks.}
Complex-valued~\cite{trabelsi2018complex} and quaternion networks~\cite{gaudet2018quaternion,zhu2018qcnn,parcollet2018qcnn,parcollet2019qrnn}
replace real weights with hypercomplex ones and the product with the Hamilton product. This shares
weights across components, cutting parameters by up to the algebra's dimension, and encodes
inter-component relations as a structural prior. Parameterised hypercomplex multiplication (PHM)
learns the multiplication rule itself for any dimension $n$~\cite{zhang2021phm}, and was extended to
convolutions~\cite{grassucci2022phnn}. Higher algebras (octonions and sedenions) have been applied to
spatiotemporal forecasting and segmentation, matching or beating real networks at a fraction of the
parameters~\cite{bojesomo2020sedenion,bojesomo2023octonion}. A unified treatment of these layers for
spatiotemporal data~\cite{bojesomo2024hypercomplex} provides the layers used in this work. Existing
hypercomplex models mostly assign the \emph{channels} of one input to the algebra's components (for
example the colour channels of an image). To our knowledge, the components have not been used to
carry distinct modality views and their interactions in audio-visual fusion, nor applied to FFIA.
\todo{Check recent hypercomplex multimodal fusion papers and cite them if any are close.}

\paragraph{Efficient models and knowledge distillation.}
Efficient audio models include CNNs pretrained on AudioSet such as PANNs~\cite{kong2020panns}, and
EfficientAT, which distils transformer ensembles into MobileNet-style CNNs~\cite{schmid2023efficientat}.
Simple pooling front-ends reduce the cost of spectrogram models further~\cite{liu2023simpf}.
For video, X3D~\cite{feichtenhofer2020x3d} and MoViNets~\cite{kondratyuk2021movinets} expand or search
mobile-sized 3-D networks; S3D~\cite{xie2018s3d} separates spatial and temporal convolutions.
Knowledge distillation transfers a large teacher's soft predictions~\cite{hinton2015kd}, intermediate
representations, or the relations between samples~\cite{park2019rkd} to a smaller student~\cite{gou2021kdsurvey}.
In FFIA, U-FFIA distils single-modality teachers into its unified model~\cite{cui2024uffia}, and
HMKD~\cite{hmkd2026} distils from an EfficientAT and MoViNet teacher pair. We combine the logit,
embedding and relation levels in one objective, and ablate the choice of teacher systematically.

% ============================================================================
\section{Related Work}
\label{sec:related}

\paragraph{Fish feeding intensity assessment.}
Early FFIA relied on observers or on hand-crafted image analysis, such as background segmentation and
the tracking of collective behaviour~\cite{zhao2017appetite}. These methods were sensitive to water
surface fluctuations and computationally heavy. Deep learning moved the field to convolutional
classifiers of feeding images~\cite{zhou2019ffia}, with later refinements such as dual attention on
EfficientNet-B2~\cite{yang2021effnetb2} and multi-scale fusion on MobileNetV3~\cite{zhang2023msif}.
Single images miss the dynamics of feeding, so video methods convert RGB sequences to optical flow and
classify them with 3-D CNNs~\cite{ubina2021ffia}, at a higher computational cost. Acoustic FFIA was
introduced with the AFFIA3K dataset, where log-mel spectrograms are classified by a CNN~\cite{cui2022affia3k}.
Surveys of fish behaviour analysis~\cite{an2021behaviour,yang2021smartfish} describe both modalities as
complementary but individually fragile. The audio-visual setting was established by AV-FFIA and
U-FFIA~\cite{cui2024uffia}. That work benchmarked audio models (PANNs~\cite{kong2020panns}, AST~\cite{gong2021ast}),
video models (I3D, S3D~\cite{xie2018s3d}, ViViT), and fusion by self-attention, cross-attention and
attention bottlenecks. It then proposed a unified mixed-modality model with modality dropout and
distillation from CNN10 and S3D teachers. Later works raised the reported AV accuracy with
cross-attention fusion~\cite{acafnet2025}, dynamic gating between experts~\cite{dynagate2026},
hierarchical multimodal fusion~\cite{hmmfnet2026}, heterogeneous distillation~\cite{hmkd2026} and
lightweight dual-modal designs~\cite{ldmf2026}. Each of them uses its own protocol: the official split,
cross-validation or a filtered subset. We report results under several protocols, including a
session-disjoint split, so that comparisons with these works are explicit.

\paragraph{Audio-visual fusion.}
Multimodal fusion is commonly divided into early, late and intermediate fusion~\cite{baltrusaitis2019mmml}.
Transformer fusion concatenates the token sequences, or uses cross-modal attention to let one modality
query the other. Attention bottlenecks restrict cross-modal flow to a few shared tokens~\cite{nagrani2021mbt},
which lowers cost and forces the model to condense information. Audio can also be used to make video
processing cheaper by previewing which frames matter~\cite{gao2020listen}. Self-supervised
audio-visual pretraining, such as contrastive masked autoencoders~\cite{gong2023cavmae}, learns joint
representations from large unlabelled corpora. Modality dropout~\cite{neverova2016moddrop} makes a fused model
robust to missing inputs and is used by U-FFIA to train a single model for audio, video or both. HCF
keeps the shared bottleneck tokens of~\cite{nagrani2021mbt}, but differs in two ways: the bottleneck
reads both modalities \emph{in parallel}, and their views and interactions are combined by an
algebraic product rather than by sequential attention or concatenation. We further attach HCF at any
subset of depths, with feedback into both streams.

\paragraph{Hypercomplex neural networks.}
Complex-valued~\cite{trabelsi2018complex} and quaternion networks~\cite{gaudet2018quaternion,zhu2018qcnn,parcollet2018qcnn,parcollet2019qrnn}
replace real weights with hypercomplex ones and the product with the Hamilton product. This shares
weights across components, cutting parameters by up to the algebra's dimension, and encodes
inter-component relations as a structural prior. Parameterised hypercomplex multiplication (PHM)
learns the multiplication rule itself for any dimension $n$~\cite{zhang2021phm}, and was extended to
convolutions~\cite{grassucci2022phnn}. Higher algebras (octonions and sedenions) have been applied to
spatiotemporal forecasting and segmentation, matching or beating real networks at a fraction of the
parameters~\cite{bojesomo2020sedenion,bojesomo2023octonion}. A unified treatment of these layers for
spatiotemporal data~\cite{bojesomo2024hypercomplex} provides the layers used in this work. Existing
hypercomplex models mostly assign the \emph{channels} of one input to the algebra's components (for
example the colour channels of an image). To our knowledge, the components have not been used to
carry distinct modality views and their interactions in audio-visual fusion, nor applied to FFIA.
\todo{Check recent hypercomplex multimodal fusion papers and cite them if any are close.}

\paragraph{Efficient models and knowledge distillation.}
Efficient audio models include CNNs pretrained on AudioSet such as PANNs~\cite{kong2020panns}, and
EfficientAT, which distils transformer ensembles into MobileNet-style CNNs~\cite{schmid2023efficientat}.
Simple pooling front-ends reduce the cost of spectrogram models further~\cite{liu2023simpf}.
For video, X3D~\cite{feichtenhofer2020x3d} and MoViNets~\cite{kondratyuk2021movinets} expand or search
mobile-sized 3-D networks; S3D~\cite{xie2018s3d} separates spatial and temporal convolutions.
Knowledge distillation transfers a large teacher's soft predictions~\cite{hinton2015kd}, intermediate
representations, or the relations between samples~\cite{park2019rkd} to a smaller student~\cite{gou2021kdsurvey}.
In FFIA, U-FFIA distils single-modality teachers into its unified model~\cite{cui2024uffia}, and
HMKD~\cite{hmkd2026} distils from an EfficientAT and MoViNet teacher pair. We combine the logit,
embedding and relation levels in one objective, and ablate the choice of teacher systematically.

% ============================================================================
\section{Related Work}
\label{sec:related}

\paragraph{Fish feeding intensity assessment.}
Early FFIA relied on observers or on hand-crafted image analysis, such as background segmentation and
the tracking of collective behaviour~\cite{zhao2017appetite}. These methods were sensitive to water
surface fluctuations and computationally heavy. Deep learning moved the field to convolutional
classifiers of feeding images~\cite{zhou2019ffia}, with later refinements such as dual attention on
EfficientNet-B2~\cite{yang2021effnetb2} and multi-scale fusion on MobileNetV3~\cite{zhang2023msif}.
Single images miss the dynamics of feeding, so video methods convert RGB sequences to optical flow and
classify them with 3-D CNNs~\cite{ubina2021ffia}, at a higher computational cost. Acoustic FFIA was
introduced with the AFFIA3K dataset, where log-mel spectrograms are classified by a CNN~\cite{cui2022affia3k}.
Surveys of fish behaviour analysis~\cite{an2021behaviour,yang2021smartfish} describe both modalities as
complementary but individually fragile. The audio-visual setting was established by AV-FFIA and
U-FFIA~\cite{cui2024uffia}. That work benchmarked audio models (PANNs~\cite{kong2020panns}, AST~\cite{gong2021ast}),
video models (I3D, S3D~\cite{xie2018s3d}, ViViT), and fusion by self-attention, cross-attention and
attention bottlenecks. It then proposed a unified mixed-modality model with modality dropout and
distillation from CNN10 and S3D teachers. Later works raised the reported AV accuracy with
cross-attention fusion~\cite{acafnet2025}, dynamic gating between experts~\cite{dynagate2026},
hierarchical multimodal fusion~\cite{hmmfnet2026}, heterogeneous distillation~\cite{hmkd2026} and
lightweight dual-modal designs~\cite{ldmf2026}. Each of them uses its own protocol: the official split,
cross-validation or a filtered subset. We report results under several protocols, including a
session-disjoint split, so that comparisons with these works are explicit.

\paragraph{Audio-visual fusion.}
Multimodal fusion is commonly divided into early, late and intermediate fusion~\cite{baltrusaitis2019mmml}.
Transformer fusion concatenates the token sequences, or uses cross-modal attention to let one modality
query the other. Attention bottlenecks restrict cross-modal flow to a few shared tokens~\cite{nagrani2021mbt},
which lowers cost and forces the model to condense information. Audio can also be used to make video
processing cheaper by previewing which frames matter~\cite{gao2020listen}. Self-supervised
audio-visual pretraining, such as contrastive masked autoencoders~\cite{gong2023cavmae}, learns joint
representations from large unlabelled corpora. Modality dropout~\cite{neverova2016moddrop} makes a fused model
robust to missing inputs and is used by U-FFIA to train a single model for audio, video or both. HCF
keeps the shared bottleneck tokens of~\cite{nagrani2021mbt}, but differs in two ways: the bottleneck
reads both modalities \emph{in parallel}, and their views and interactions are combined by an
algebraic product rather than by sequential attention or concatenation. We further attach HCF at any
subset of depths, with feedback into both streams.

\paragraph{Hypercomplex neural networks.}
Complex-valued~\cite{trabelsi2018complex} and quaternion networks~\cite{gaudet2018quaternion,zhu2018qcnn,parcollet2018qcnn,parcollet2019qrnn}
replace real weights with hypercomplex ones and the product with the Hamilton product. This shares
weights across components, cutting parameters by up to the algebra's dimension, and encodes
inter-component relations as a structural prior. Parameterised hypercomplex multiplication (PHM)
learns the multiplication rule itself for any dimension $n$~\cite{zhang2021phm}, and was extended to
convolutions~\cite{grassucci2022phnn}. Higher algebras (octonions and sedenions) have been applied to
spatiotemporal forecasting and segmentation, matching or beating real networks at a fraction of the
parameters~\cite{bojesomo2020sedenion,bojesomo2023octonion}. A unified treatment of these layers for
spatiotemporal data~\cite{bojesomo2024hypercomplex} provides the layers used in this work. Existing
hypercomplex models mostly assign the \emph{channels} of one input to the algebra's components (for
example the colour channels of an image). To our knowledge, the components have not been used to
carry distinct modality views and their interactions in audio-visual fusion, nor applied to FFIA.
\todo{Check recent hypercomplex multimodal fusion papers and cite them if any are close.}

\paragraph{Efficient models and knowledge distillation.}
Efficient audio models include CNNs pretrained on AudioSet such as PANNs~\cite{kong2020panns}, and
EfficientAT, which distils transformer ensembles into MobileNet-style CNNs~\cite{schmid2023efficientat}.
Simple pooling front-ends reduce the cost of spectrogram models further~\cite{liu2023simpf}.
For video, X3D~\cite{feichtenhofer2020x3d} and MoViNets~\cite{kondratyuk2021movinets} expand or search
mobile-sized 3-D networks; S3D~\cite{xie2018s3d} separates spatial and temporal convolutions.
Knowledge distillation transfers a large teacher's soft predictions~\cite{hinton2015kd}, intermediate
representations, or the relations between samples~\cite{park2019rkd} to a smaller student~\cite{gou2021kdsurvey}.
In FFIA, U-FFIA distils single-modality teachers into its unified model~\cite{cui2024uffia}, and
HMKD~\cite{hmkd2026} distils from an EfficientAT and MoViNet teacher pair. We combine the logit,
embedding and relation levels in one objective, and ablate the choice of teacher systematically.

% ============================================================================
\section{Related Work}
\label{sec:related}

\paragraph{Fish feeding intensity assessment.}
Early FFIA relied on observers or on hand-crafted image analysis, such as background segmentation and
the tracking of collective behaviour~\cite{zhao2017appetite}. These methods were sensitive to water
surface fluctuations and computationally heavy. Deep learning moved the field to convolutional
classifiers of feeding images~\cite{zhou2019ffia}, with later refinements such as dual attention on
EfficientNet-B2~\cite{yang2021effnetb2} and multi-scale fusion on MobileNetV3~\cite{zhang2023msif}.
Single images miss the dynamics of feeding, so video methods convert RGB sequences to optical flow and
classify them with 3-D CNNs~\cite{ubina2021ffia}, at a higher computational cost. Acoustic FFIA was
introduced with the AFFIA3K dataset, where log-mel spectrograms are classified by a CNN~\cite{cui2022affia3k}.
Surveys of fish behaviour analysis~\cite{an2021behaviour,yang2021smartfish} describe both modalities as
complementary but individually fragile. The audio-visual setting was established by AV-FFIA and
U-FFIA~\cite{cui2024uffia}. That work benchmarked audio models (PANNs~\cite{kong2020panns}, AST~\cite{gong2021ast}),
video models (I3D, S3D~\cite{xie2018s3d}, ViViT), and fusion by self-attention, cross-attention and
attention bottlenecks. It then proposed a unified mixed-modality model with modality dropout and
distillation from CNN10 and S3D teachers. Later works raised the reported AV accuracy with
cross-attention fusion~\cite{acafnet2025}, dynamic gating between experts~\cite{dynagate2026},
hierarchical multimodal fusion~\cite{hmmfnet2026}, heterogeneous distillation~\cite{hmkd2026} and
lightweight dual-modal designs~\cite{ldmf2026}. Each of them uses its own protocol: the official split,
cross-validation or a filtered subset. We report results under several protocols, including a
session-disjoint split, so that comparisons with these works are explicit.

\paragraph{Audio-visual fusion.}
Multimodal fusion is commonly divided into early, late and intermediate fusion~\cite{baltrusaitis2019mmml}.
Transformer fusion concatenates the token sequences, or uses cross-modal attention to let one modality
query the other. Attention bottlenecks restrict cross-modal flow to a few shared tokens~\cite{nagrani2021mbt},
which lowers cost and forces the model to condense information. Audio can also be used to make video
processing cheaper by previewing which frames matter~\cite{gao2020listen}. Self-supervised
audio-visual pretraining, such as contrastive masked autoencoders~\cite{gong2023cavmae}, learns joint
representations from large unlabelled corpora. Modality dropout~\cite{neverova2016moddrop} makes a fused model
robust to missing inputs and is used by U-FFIA to train a single model for audio, video or both. HCF
keeps the shared bottleneck tokens of~\cite{nagrani2021mbt}, but differs in two ways: the bottleneck
reads both modalities \emph{in parallel}, and their views and interactions are combined by an
algebraic product rather than by sequential attention or concatenation. We further attach HCF at any
subset of depths, with feedback into both streams.

\paragraph{Hypercomplex neural networks.}
Complex-valued~\cite{trabelsi2018complex} and quaternion networks~\cite{gaudet2018quaternion,zhu2018qcnn,parcollet2018qcnn,parcollet2019qrnn}
replace real weights with hypercomplex ones and the product with the Hamilton product. This shares
weights across components, cutting parameters by up to the algebra's dimension, and encodes
inter-component relations as a structural prior. Parameterised hypercomplex multiplication (PHM)
learns the multiplication rule itself for any dimension $n$~\cite{zhang2021phm}, and was extended to
convolutions~\cite{grassucci2022phnn}. Higher algebras (octonions and sedenions) have been applied to
spatiotemporal forecasting and segmentation, matching or beating real networks at a fraction of the
parameters~\cite{bojesomo2020sedenion,bojesomo2023octonion}. A unified treatment of these layers for
spatiotemporal data~\cite{bojesomo2024hypercomplex} provides the layers used in this work. Existing
hypercomplex models mostly assign the \emph{channels} of one input to the algebra's components (for
example the colour channels of an image). To our knowledge, the components have not been used to
carry distinct modality views and their interactions in audio-visual fusion, nor applied to FFIA.
\todo{Check recent hypercomplex multimodal fusion papers and cite them if any are close.}

\paragraph{Efficient models and knowledge distillation.}
Efficient audio models include CNNs pretrained on AudioSet such as PANNs~\cite{kong2020panns}, and
EfficientAT, which distils transformer ensembles into MobileNet-style CNNs~\cite{schmid2023efficientat}.
Simple pooling front-ends reduce the cost of spectrogram models further~\cite{liu2023simpf}.
For video, X3D~\cite{feichtenhofer2020x3d} and MoViNets~\cite{kondratyuk2021movinets} expand or search
mobile-sized 3-D networks; S3D~\cite{xie2018s3d} separates spatial and temporal convolutions.
Knowledge distillation transfers a large teacher's soft predictions~\cite{hinton2015kd}, intermediate
representations, or the relations between samples~\cite{park2019rkd} to a smaller student~\cite{gou2021kdsurvey}.
In FFIA, U-FFIA distils single-modality teachers into its unified model~\cite{cui2024uffia}, and
HMKD~\cite{hmkd2026} distils from an EfficientAT and MoViNet teacher pair. We combine the logit,
embedding and relation levels in one objective, and ablate the choice of teacher systematically.
